%% =============================================================================
%% sm_c2_propriety.tex — SM-C, Subsection C.2: Pseudo-Posterior Propriety
%% Body content only (\subsection level); parent structure in sm_c.tex
%% Primary source: dev/log/06_research-notes-survey-design-pseudo-posterior.qmd
%%   lines 741--781 (Definition 3.9, Theorem 3.1, Remarks 3.8--3.9)
%% =============================================================================

\subsection{Pseudo-posterior propriety}
\label{smc:propriety}

We now provide the complete proof of~\cref{thm:propriety}, which was
stated with a sketch in~\cref{sec:survey}.

\begin{proof}[Proof of \textup{\cref{thm:propriety}}]
  We verify that the pseudo-posterior
  \[
    \pi^{(w)}(\btheta \mid \by) \propto
    \exp\bigl[\ell^{(w)}(\btheta)\bigr]\,\pi(\btheta)
  \]
  is integrable over the parameter space.

  \textup{\textbf{Step 1} (Boundedness of the pseudo-likelihood).}
  For each observation $i$, the HBB probability satisfies
  $f_{\HBB}(y_i \mid \btheta) \in (0, 1]$ for all $\btheta$ in the
  interior of the parameter space.  Since $\tilde{w}_i > 0$, we have
  $f_i^{\tilde{w}_i} \le 1$ whenever $f_i \le 1$.  Therefore
  \[
    \exp\bigl[\ell^{(w)}(\btheta)\bigr]
    = \prod_{i=1}^{N} f_i^{\tilde{w}_i}
    \le 1
    \qquad \text{for all } \btheta.
  \]

  \textup{\textbf{Step 2} (Properness of the prior).}
  Each component of the prior $\pi(\btheta)$ is a proper distribution:
  \begin{itemize}[nosep]
    \item $\balpha, \bbeta \sim \Norm(\mathbf{0},\; 2^2\,\mathbf{I}_P)$
      \citep{GelmanEtAl2008}: proper with finite variance;
    \item $\boldsymbol{\tau} \sim \HalfNormal(0, 1)$ componentwise:
      proper on $\R_{>0}$;
    \item $\mathbf{R}_\varepsilon \sim \LKJ(2)$
      \citep{LewandowskiEtAl2009}: proper over the space of
      $2q \times 2q$ correlation matrices;
    \item $\log\kappa \sim \Norm(2,\; 1.5^2)$: proper.
  \end{itemize}
  The joint prior is a product of proper distributions, hence proper:
  $\int \pi(\btheta)\,d\btheta = 1$.

  \textup{\textbf{Step 3} (Integrability).}
  Combining Steps~1 and~2:
  \[
    \int \exp\bigl[\ell^{(w)}(\btheta)\bigr]\,\pi(\btheta)\,d\btheta
    \;\le\;
    \int 1 \cdot \pi(\btheta)\,d\btheta
    = 1 < \infty.
  \]
  Moreover, $\exp[\ell^{(w)}(\btheta)] > 0$ on the interior of the
  parameter space (since all probability terms $f_i$ are strictly
  positive there), so the pseudo-posterior is non-degenerate.
\end{proof}

\begin{remark}[Weight-invariance of propriety]
\label{smc:rem-weight-invariance}
  The proof above does not depend on the specific values or
  normalization of the weights---only on their positivity and
  finiteness.  Propriety therefore holds regardless of whether the
  weights are normalized to sum to~$N$, to sum to the population size
  $\hat{M}$, or under any other positive scaling.  However, the
  \emph{shape} and \emph{concentration} of the pseudo-posterior depend
  critically on the normalization convention, as discussed in
  \cref{smc:rem-normalization}: the sum-to-$\hat{M}$ convention
  produces extreme overconcentration, while the sum-to-$N$ convention
  preserves the information content of the actual sample.
\end{remark}

\begin{remark}[Contrast with improper priors]
\label{smc:rem-improper-priors}
  With improper (flat) priors on the regression coefficients,
  propriety would require the pseudo-likelihood to be integrable on
  its own---that is, $\int \exp[\ell^{(w)}(\btheta)]\,d\btheta < \infty$.
  In this setting the effective sample size (not the nominal sample
  size) governs integrability.  With $\ESS \approx 1{,}797$ and
  $\dim(\btheta) \le 606$, integrability would likely hold, but the
  formal verification is more delicate: it requires confirming that the
  weighted Hessian has full rank at the pseudo-MLE for every parameter
  block.  The use of proper priors sidesteps this issue entirely,
  ensuring propriety for \emph{any} positive weight configuration
  without conditions on the data.
\end{remark}
